% 第 59 章　规范对称性与相互作用（挑战篇）
\chapter{规范对称性与相互作用（挑战篇）}

第 57 章弄清了什么叫对称：做某种改变之后，规律不变。第 58 章看到了对称性的惊人产出：每一种连续对称都对应一条守恒律。这一章要把问题拧到最紧的一档——如果那个“改变”根本什么都没改，只是换了一套描述方式呢？比如把海拔的零点从海平面挪到你家地板，把波函数的相位整体转一个角度，物理规律该不该在乎？答案是“不该在乎”，这听上去像一句废话。可二十世纪的物理学家发现，把这句废话当成一条严格的设计要求，竟然能把电磁力、弱力和强力这三种相互作用的骨架“要求”出来。这条线索叫规范对称性，它是现代粒子物理最深的引擎，也是本章要攀的山。本章是挑战篇，数学会多一些，但每一步仍然先给直觉，再给符号。

\step{反常识开场}

\begin{mainline}
先看一个能在窗边观察的场面。高压输电线上，小鸟一只挨一只地站着，悠然自得。那根导线的电压以万伏计，有的线路高达五十万伏；而家里墙上的插座只有 220 伏，却年年伤人。电死的到底是什么？“电压高”这三个字显然解释不通：五十万伏的电线对鸟秋毫无犯，220 伏的插座对人足以致命。看来“某一点的电压是多少”这个问题本身就问错了——真正要紧的是别的东西。

再看一个只能在实验室里上演、却被反复验证过的怪事。把电子一个一个地射向双缝，屏上慢慢积累出干涉条纹——这是量子篇里熟悉的老朋友。现在，在两条缝的正后方、电子必经区域的中间，藏进一根极细的螺线管。这种螺线管有个好脾气：通上电流后，磁场被严严实实地锁在自己的肚子里，管外的磁场处处为零。电子从管子两侧飞过，它们的飞行路径上连一丝磁场都检测不到。按理说，什么都不会改变。可是实验者合上开关的瞬间，屏上的干涉条纹移动了位置；断开开关，条纹移回去。电子从头到尾没有进过磁场，没有受到任何力的作用，它凭什么知道螺线管开了？

这个实验叫阿哈罗诺夫—玻姆效应，1959 年从理论上被预言，此后被一代代实验反复确认。1986 年，外村彰领导的日本团队做了一个近乎无懈可击的版本：他们用环形磁体代替螺线管，让磁场沿着环闭合，一滴都不漏到外面，还在磁体外面包了超导屏蔽层，彻底杜绝电子“漏进”磁场区域的可能。条纹照样移动。一边是小鸟脚下的五十万伏“不存在”，一边是电子路径上不存在的磁场“起了作用”。这两个反常凑在一起，指向同一个谜：物理学里到底哪些量是真实的，哪些只是我们描述世界的“记账习惯”？这个谜的完整答案，就是本章。
\end{mainline}

\step{先猜一猜}

\begin{mainline}
在揭晓之前，先把你自己的判断写下来。下面四句话，凭直觉判断对错：

第一，电压表测量的是“某一点的电压”。

第二，小鸟在高压线上安然无恙，是因为鸟爪的角质层绝缘特别好，电流根本进不了鸟的身体。

第三，磁场为零的地方，任何物体都不可能感知到一块磁铁的存在——力都无从谈起，影响更无从谈起。

第四，假如全世界开会决定，把海拔零点从海平面改到珠峰大本营，所有山的高度都要改写，那么洪水预警、飞机高度表、水坝设计全都要推倒重算。

请先写下你的答案。读完本章再回头对照：第一句话有个藏得很深的语病；第二句话的方向对了一半，但真正的原因与绝缘关系不大；第三句话被本章最核心的实验直接推翻；第四句话里，哪些计算真的不用重算、哪些概念根本不依赖零点，恰恰是本章前半段的主题。
\end{mainline}

\step{动手实验}

\begin{experiment}[换个零点，水柱答应吗]
找一个大饮料瓶，在侧面不同高度扎两个小孔，孔径差不多。在瓶底贴一条胶带当“零刻度线”，往上量出两个孔的高度，记作 \(h_1\) 和 \(h_2\)。把瓶子放在地上，装满水，水面上方敞开，观察两股水柱的射程：下面的孔喷得远，上面的孔喷得近。这符合直觉——越深处压强越大。

现在做本章的关键动作：把整个瓶子端到一米高的椅子上，再把“零刻度线”撕下来贴到桌面上重新量。两个孔的“海拔”全都变了：有的读者量出来每个孔的高度多了几十厘米，有的读者从桌面量，数值干脆变成了负数。可你盯着水柱看——射程一丝一毫都没有变。

原因在第 12 章讲过：液面下深度为 \(d\) 处的压强比液面高出 \(\rho g d\)，决定水柱速度的是这个压强差，而压强差只认“到水面的竖直距离”，不认你贴在瓶子上的任何一条零刻度线。把零点挪到任何地方，每个高度的读数都同时加上或减去同一个数，差值纹丝不动。物理规律关心的从来不是高度本身，而是高度的差。刚才你亲手验证的这件事，有一个正式的名字：势能零点的规范自由。
\end{experiment}

\begin{mainline}
还有一个零成本的版本。打开手机里的气压计或海拔应用，在楼下记一个读数，爬到楼顶再记一个。你会发现两件事：一是不同品牌的手机、不同的应用给出的绝对海拔可能差出十几米——因为它们的“零点”换算各有各的来源；二是两部手机测出的高度差却几乎完全一致，和楼层实际高度对得上。绝对读数是约定，差值才是物理。请记住这个句型，本章会把它用到量子力学的相位上：绝对相位是约定，相位差才是物理。
\end{mainline}

\step{旧模型遇到麻烦}

\begin{mainline}
现在把旧观念摊开，看看它究竟在哪里开始漏水。麻烦一共有三处。

第一处麻烦：“电势只是记账工具”这个说法，被开场那个电子实验顶了回来。第 20 章引入电势时，我们说得很坦白：电场是真实的，它能推得动电荷；电势是为了省事发明的地形图，是电场一路积分出来的账簿。按这个观念，电子飞行路径上电场和磁场都为零，电子就什么效应都不该有。可是条纹移动了。账簿自己站起来影响了现实，“工具”这个词罩不住它了。

第二处麻烦来自反方向：“电势是真实物理量”这个说法同样站不住。请你设计任何一个测量“某一点绝对电势”的实验——注意，不是电势差，是电势本身。你会发现根本设计不出来。电压表有两个探头，它读的永远是两点之差；鸟站在电线上平安无事，因为它全身只接触一个电势；开场那瓶水也只认高度差。整个物理学里没有任何仪器、任何实验能读出“这一点比无穷远高多少伏”这个数的绝对意义，除非你先把“无穷远是零”这条人为约定写进测量方案。一个原则上测不到的量，有资格叫“真实”吗？

第三处麻烦最深，藏在前两处的夹缝里。量子篇（第 43 章和第 44 章）讲过，量子态由一个波函数描述，波函数带着相位，相位可以画成钟面上的一根指针。有两件事那时就埋下了伏笔。其一，把全宇宙所有波函数的相位指针同时转同一个角度，任何测量结果都不变——干涉只认相位差，不认绝对相位。这和“海拔零点随便挪”完全同构。其二，不同路径之间的相对相位是命根子，双缝干涉的亮纹暗纹全由它决定。现在把两件事撞在一起：既然绝对相位无所谓，我能不能让每个地点各用各的相位零点——北京按北京的钟面，上海按上海的钟面？如果允许，一个“自由粒子”的波函数从北京传到上海，它的相位变化就没法定义了，薛定谔方程当场散架；如果不允许，凭什么？物理规律凭什么要求全世界共用同一个相位零点？旧模型在这个问题上哑口无言，而这正是新概念的入口。
\end{mainline}

\step{发明新概念}

\begin{mainline}
先把两种“随便改”严格分开，这是本章第一个分水岭。

第一种叫全局规范自由：全世界一起改。所有海拔同时减去十米，所有波函数的相位指针同时转三十度。这种改变连描述的“纹理”都没动，只是整体平移了一下标签，物理当然不变。第 58 章的诺特定理还在这里等着收钱：波函数在整体相位转动下不变，这条对称性对应的守恒量，正是电荷。电荷守恒不是什么神秘禁令，它是“绝对相位不可测”的另一张面孔。

第二种叫局域规范自由：每个地点各改各的。北京的相位指针转三十度，上海的转负十五度，杭州的不动。这就不再是换标签那么简单了——它改变了相邻地点之间的“相对关系”，而相对相位是干涉的命根子，是真实的物理。要求物理规律在这种逐点任意重选零点的操作下依然成立，是一条苛刻到近乎无理的要求。可 1929 年外尔（更早他在 1918 年试过另一个版本）、1954 年杨振宁和米尔斯沿着这条路走下去，发现了一件震惊物理学界的事：如果坚持这条要求，自由粒子的量子理论是不合格的，必须往方程里加一个新的场；这个场在每一点的取值，恰好扮演“相邻地点相位零点之间换算规则”的角色；而它与带电粒子的耦合方式被这条要求几乎唯一地钉死——它就是电磁场。电磁相互作用不是被“发现”之后再找解释，而是可以被“相位零点任选”这条设计原则反推出来。这条思路叫规范原理。

它像什么？它像时区制度。每个城市按自己的太阳定本地时间，这就是“每点各选各的零点”；而航空公司的时刻表必须附带换算规则，告诉你飞机从上海飞到法兰克福，到达时刻减去起飞时刻再减去时差，才是真正的飞行时长。电磁场中的矢势 \(\mathbf A\)，就是宇宙级别的“时差换算表”。它不像什么？它不像一只挂在天上的标准钟，强迫所有城市对表——恰恰相反，它承认每个地点有选零点的自由，然后把“比较两地”的责任揽到自己身上。

但有一个警告必须立刻插上，这是本章的护栏：规范对称性与第 57 章讲的旋转、平移对称不是同一种东西。旋转一个物体，你得到的是另一个物理状态——它确实朝另一个方向了；而做一次规范变换，你没有得到任何新状态，只是给同一个状态换了一套描述。所以规范“对称”更诚实叫法是描述的冗余。可是阿哈罗诺夫—玻姆效应警告我们，不能反过来滑到另一个极端，说矢势 \(\mathbf A\) 纯属虚构。真正的结论是：\(\mathbf A\) 的逐点取值是冗余的（可以被规范变换改来改去），而它沿闭合回路的积分——也就是回路所围的磁通——是真实的、可测的、改不掉的。冗余与真实共存在一个对象里，这正是规范理论微妙而有力的地方。

把这条线索推到底，就得到了一张现代物理的概念地图。电磁相互作用对应的相位自由是“一根指针在圆盘上转”，数学上叫 \(U(1)\) 群，它的传播者是光子，没有质量，力程无限。弱相互作用对应的是一种双分量量子系统（弱同位旋）内部的“相位”自由，数学上叫 \(SU(2)\)，它的传播者有三个：\(W^{+}\)、\(W^{-}\) 和 \(Z\)。强相互作用对应夸克三种“颜色”之间的重排自由，数学上叫 \(SU(3)\)，传播者是八种胶子。三种力、三种“每点任选约定”的自由、三套传播者，结构上是同一个模子刻出来的。这张地图不是空想：\(W\) 和 \(Z\) 粒子于 1983 年在欧洲核子中心的对撞机里被直接找到，质量与弱电理论的预言吻合；胶子虽因禁闭无法单独现身，但它在高能对撞中留下的“三喷注”信号在 1979 年就被确认。地图上最大的那块空白——\(W\) 和 \(Z\) 为什么有质量、弱力为什么短程——正是第 60 章对称性破缺要填的。

细心的读者会问：引力呢？四种基本相互作用里，为什么这张地图只有三家？诚实的回答是：引力的故事不一样。第 41 章会讲到，引力最深刻的描述不是“一种力”，而是时空本身的弯曲，它的“对称性”牵涉的是时空坐标而非内部相位。把引力也纳入规范框架的努力一直在进行，并取得了部分成功，但把引力量子化的完整理论至今仍是物理学最大的未完工工地。这张地图留出的那把空椅子，会在第 61 章的全书地图里再摆出来一次。
\end{mainline}

\step{一句话模型}

\begin{oneline}
物理预测不许依赖描述中的人为选择：电势的零点、波函数相位的零点都可以任意挪；要求“每个地点各自任挪”之后规律仍然成立，方程里就必须存在一个负责在相邻地点之间做换算的场——电磁场（以及它的弱、强表亲）就是这样被“必须不依赖约定”这条原则逼出来的；场的逐点取值是约定，而回路通量与场强才是甩不掉的物理。
\end{oneline}

\step{数学翻译器}

\begin{mainline}
先把相位画出来。一个量子态在某地的相位，就是钟面上的一根指针；指针本身指向几点不重要，两根指针的夹角才重要。双缝干涉里，屏上某一点是亮是暗，取决于“沿缝一路径到达的指针”和“沿缝二路径到达的指针”之间的夹角：夹角为零（指向相同）就是亮纹，夹角为半个圆周（指向相反）就是暗纹。
\end{mainline}

\begin{center}
\begin{tikzpicture}[scale=1.0,>=Stealth]
  % 钟面一
  \draw[thick] (0,0) circle (0.9);
  \draw[thick,->] (0,0) -- (35:0.85);
  \node at (0,-1.3) {地点甲的相位指针};
  % 钟面二
  \draw[thick] (4.5,0) circle (0.9);
  \draw[thick,->] (4.5,0) -- (110:0.85);
  \node at (4.5,-1.3) {地点乙的相位指针};
  % 换算表
  \draw[decorate,decoration={snake,amplitude=0.4mm}] (0.95,0.3) .. controls (2.2,1.4) .. (3.55,0.3);
  \node[above] at (2.25,1.15) {矢势 \(\mathbf A\)：相邻零点的换算规则};
\end{tikzpicture}
\end{center}

\begin{mainline}
现在写出本章的核心公式。带电 \(q\) 的粒子沿一条路径从甲地走到乙地，矢势 \(\mathbf A\) 会给它的相位指针一个额外的转动，转过的角度是
\[
\Delta\varphi=\frac{q}{\hbar}\int_{\text{路径}}\mathbf A\cdot d\mathbf l
\]
其中 \(\hbar\) 是约化普朗克常数，积分表示沿路径把 \(\mathbf A\) 的顺路分量一小段一小段累加起来。这个公式描述什么？描述电磁环境如何拨动量子相位的指针。为什么这样写？因为它是唯一能同时满足三条硬约束的写法：路径上磁场为零时电场也为零的普通情形下它给出零额外相位；它在“换一套零点约定”（规范变换）下对所有闭合比较给出相同结论；它与第 20 章电势差的故事无缝衔接。何时成立？带电粒子在磁场可以逐点定义的区域里运动。何时失效？马上就会看到——当几条路径围成回路时，即使路径上处处无磁场，它也能给出非零的答案。

把两条路径拼成一个回路，从源点分两路出发、在屏上会合，相对相位就是沿回路的积分。数学里有一条定理（斯托克斯定理）保证：矢势沿回路的积分等于穿过回路所围面积的磁通 \(\Phi\)。于是
\[
\Delta\varphi_{\text{相对}}=\frac{q\Phi}{\hbar}
\]
屏上的条纹整体平移，平移量占一个条纹间距的比例是 \(\Delta\varphi/2\pi=\Phi/\Phi_0\)，其中
\[
\Phi_0=\frac{h}{e}\approx4.14\times10^{-15}\ \text{韦伯}
\]
叫磁通量子——电子的电荷让“一圈相位转满”对应的磁通成为一个自然单位。

按老规矩，立刻给公式做体检。零检查：磁通为零，相对相位为零，条纹不动，与没有螺线管时一致。极大检查：磁通很大时相位很大，但相位转满整圈后指针回到原处，所以条纹的移动永远是周期性的——每多一个磁通量子，图样恰好平移一个条纹间距，周而复始。方向反转检查：螺线管电流反向，磁通变号，相位变号，条纹反向平移。极小检查：磁通等于半个磁通量子时，相位差半个圆周，亮纹暗纹恰好互换位置。每一条都能直接拿去实验室核对，这正是这个公式令人敬畏的地方。
\end{mainline}

\begin{center}
\begin{tikzpicture}[scale=1.0,>=Stealth]
  % 电子源
  \fill (0,1) circle (0.08);
  \node[left] at (-0.1,1) {电子源};
  % 双缝挡板
  \draw[thick] (1.5,0) -- (1.5,0.7);
  \draw[thick] (1.5,0.85) -- (1.5,1.15);
  \draw[thick] (1.5,1.3) -- (1.5,2);
  % 螺线管
  \draw[thick,fill=gray!20] (2.6,1) circle (0.22);
  \node[below] at (2.6,0.7) {螺线管（磁场锁在内部）};
  % 路径
  \draw[thick] (0,1) -- (1.5,0.95) -- (4.5,0.55);
  \draw[thick] (0,1) -- (1.5,1.05) -- (4.5,1.45);
  \node[above] at (1.1,1.55) {路径二};
  \node[below] at (1.1,0.42) {路径一};
  % 屏
  \draw[thick] (4.5,0.2) -- (4.5,1.8);
  \node[below] at (4.5,0.2) {探测屏};
  % 条纹
  \foreach \y in {0.45,0.7,0.95,1.2,1.45} \draw (4.62,\y) -- (5.1,\y);
  \node[right] at (5.15,1.0) {条纹};
\end{tikzpicture}
\end{center}

\begin{mainline}
做一笔带真实数据的估算。取一根半径一微米（头发丝直径的七十分之一）的微型螺线管，内部磁场只有 \(0.01\) 特斯拉——一块普通小磁铁表面磁场的十分之一左右，轻松可达。它锁住的磁通是 \(\Phi=B\pi r^{2}\approx0.01\times3.14\times10^{-12}\approx3\times10^{-14}\) 韦伯。除以磁通量子：\(\Phi/\Phi_0\approx3\times10^{-14}/4.14\times10^{-15}\approx7.5\)。也就是说，合上开关，干涉图样整体平移约七个半条纹间距——在电子显微镜下这是排山倒海的变化，绝不可能看漏。而与此同时，电子飞行路径上的磁场测量值是零，逐点都是零。力是零，效应却不为零，这就是“磁场为零的地方也能感知磁铁”的精确含义：感知的不是磁场，是路径围不住的磁通。
\end{mainline}

\begin{mathbridge}
\textbf{一、相位为什么写成指针的乘法。} 量子态的相位因子写作 \(e^{i\theta}\)，整体转动就是把波函数乘以 \(e^{i\theta_0}\)。干涉项来自两路叠加 \(\psi_1+\psi_2\) 的模方，整体因子在取模时被 \(|e^{i\theta_0}|=1\) 吃掉，所以绝对相位永不可测；而相对因子 \(e^{i(\theta_1-\theta_2)}\) 留在交叉项里，决定亮暗。这和“海拔差与零点无关”是同一条代数：可观测的永远是差。

\textbf{二、规范变换是两个人约好一起变。} 把相位零点逐点重选，数学上是 \(\psi\to e^{iq\chi(\mathbf x)/\hbar}\psi\)，其中 \(\chi(\mathbf x)\) 是任意函数。要让物理不变，矢势必须同步变换 \(\mathbf A\to\mathbf A+\nabla\chi\)。单独变任何一个，方程都会变脸；两个一起变，所有可观测量子丝不动。逐点看，\(\mathbf A\) 改了；但回路积分 \(\oint\mathbf A\cdot d\mathbf l=\Phi\) 里，\(\nabla\chi\) 沿回路积一圈回到起点，贡献恒为零——所以磁通是规范变换拿它没办法的“硬通货”。电势 \(V\) 也有对应搭档变换，本书电磁学部分用过的“零点任选”正是它的特殊情形。

\textbf{三、最小耦合：方程自己提出修改申请。} 自由粒子的薛定谔方程里，动量以导数的形式出现。逐点重选相位零点后，普通导数会把任意函数 \(\chi\) 的梯度“漏”出来，方程不再成立。补救办法几乎是被逼出来的：把动量算符换成 \(\mathbf p-q\mathbf A\)，让 \(\mathbf A\) 的变换恰好吃掉多出来的那一项。这一替换叫最小耦合，它自动给出了带电粒子与电磁场的全部相互作用——洛伦兹力、AB 相位、原子发光的选择定则，全从这一个小改动里流出来。规范原理的威力浓缩成一句话：不许依赖约定，方程就被迫长出一整门电磁学。
\end{mathbridge}

\begin{challenge}
把每个地点的相位零点想成一根圆盘上的指针，“逐点任选零点”就是在每一点做一个圆盘转动，这些转动组成的集合叫 \(U(1)\) 群。电磁学的规范理论就是 \(U(1)\) 规范理论。真正的飞跃在 1954 年：杨振宁和米尔斯问，如果每个地点的自由度不是一根指针，而是一个双分量量子系统（比如“质子—中子”这对孪生）的内部转动呢？这种转动组成的群叫 \(SU(2)\)。圆盘转动谁先谁后没有区别，而 \(SU(2)\) 的转动次序交换不得——先绕 \(x\) 轴转再绕 \(y\) 轴转，与反过来做，结果不同。这个“不可交换”逼出一个电磁学里没有的后果：传播相互作用的规范粒子自己也带着“荷”，它们彼此之间直接相互作用。强相互作用的 \(SU(3)\) 理论把这一点推到极致：胶子自己带色荷，互相纠缠，导致夸克被“禁闭”在质子内部永远拉不出来，且越靠近作用越弱（渐近自由，2004 年诺贝尔物理学奖）。电磁世界里清净的光子，只是这个家族里最不黏人的成员。

规范结构与拓扑还会联手给出预言。狄拉克在 1931 年指出：如果宇宙里存在哪怕一个磁单极子，那么波函数的相位绕它一圈必须自洽，这要求电荷与磁荷的乘积取离散的量子化值——这将一举解释“为什么所有电荷都是电子电荷的整数倍”。磁单极子至今没有被找到，但这个论证至今仍是理论物理最漂亮的范本之一：它没有用任何新实验数据，只用“相位必须说得通”这一条自洽性要求，就撬动了电荷量子化这个深层谜题。
\end{challenge}

\step{模型的边界}

\begin{mainline}
规范原理威力惊人，边界也必须划清，否则它会被用成一顶什么都能装的帽子。

第一条边界：它是设计原则，不是证明。规范原理规定的是相互作用的“形式”——必须有规范场、耦合必须取最小耦合的样子；但它不告诉你相互作用的“强度”：为什么电磁耦合常数是那个数，为什么恰好是 \(U(1)\)、\(SU(2)\)、\(SU(3)\) 这三家而不是别的群，为什么有三代夸克和轻子。这些至今是输入，不是输出。第 61 章会回到这份“我们仍然不知道”的清单。

第二条边界：经典物理学里根本没有规范原理的刚需。牛顿力学加洛伦兹力自洽地运转了两百年，从不问路面积分和相位——AB 效应是纯粹的量子效应，因为“相位”只在量子理论里出场。把规范原理讲成“经典世界早就需要的补丁”是不诚实的；它是量子理论与相对议论婚之后，嫁妆里翻出来的传家宝。

第三条边界对应典型误用。误用一：“电势、矢势是虚构的数学工具。”AB 实验直接驳回：虚构的工具移不动真实的条纹，回路磁通是硬通货。误用二：“规范对称和旋转对称一样，都是世界的对称。”不一样。旋转把物体变到另一个可区分的状态，规范变换只换描述——你不能问“规范变换之后世界变成了什么样”，因为什么都没变。把它和普通几何对称混为一谈，就会误以为它也能像诺特定理那样直接产出守恒量，这是误解。误用三：“规范理论证明了力必须存在。”它没有。它说的是：一旦接受量子理论，又要求逐点的约定自由，那么相互作用若存在就必须取规范场的形式。这是条件句，不是存在性证明。

最后按全书惯例分清三层。实验事实是：条纹随螺线管开关移动、鸟在高压线上无恙、电荷守恒在任何反应里成立、\(W\) 粒子的质量约 80 吉电子伏特——谁来测都一样。数学模型是：矢势、规范变换、\(U(1)\)、\(SU(2)\)、\(SU(3)\) 这些群、最小耦合——它们把事实组织成能算的规则。物理解释是：“相互作用是保持描述自由所付的代价”“力是约定的保镖”这类说法——它们帮助你思考，但它们是对模型的读法，不是新的事实。下一章讲对称破缺时，这个区分还会救我们一次。
\end{mainline}

\step{回头解释开场现象}

\begin{mainline}
现在回到高压线上的鸟。电流不是被“电压高”推动的，而是被电势差推动的——第 21 章的欧姆定律右边是电压差，不是电压。鸟的两只爪子抓在同一根导线上，全身处于同一个电势，爪与爪之间的电势差几乎为零，流过身体的电流也就几乎为零，与这根导线离地五十万伏还是五百万伏毫无关系。它脚下的“五十万伏”是相对于远处大地的约定差值，而大地不在它的回路里。真正的杀机出现在电势差闭合的那一刻：另一只翅膀碰到电线杆、铁塔或第二根不同相位的导线，身体瞬间变成连接两个电势的导线，惨剧发生。电压表同理：它的两个探头读的是差，所谓“测某点电压”永远是“测某点相对某个被默认的零点的差”。开场的悬念一，被“零点任选而差值为真”这条规范自由干净利落地解决——这就是猜一猜里第二句话的正确答案：与绝缘关系不大，与“是否形成回路”关系一切。

再看电子与螺线管。电子走两条路在屏上会合，两路之间围出一个回路，回路中间正好套着那根螺线管。开关断开，穿过回路的磁通为零，两条路径的相对相位就是原来的几何差，条纹停在原位；开关合上，磁通变成 \(\Phi\)，相对相位被拧动 \(e\Phi/\hbar\)，条纹整体平移 \(\Phi/\Phi_0\) 个间距。电子全程在磁场为零的区域飞行，受的力逐点为零——这点旧模型没说错；但量子世界里的舞台提示不是力，而是相位，相位认的是回路围住的磁通，磁通不因为路径上没有磁场就消失。外村彰团队 1986 年的实验把屏蔽做到极致，条纹依然按预言平移，至此“只有场才真实”的旧信条正式让位给“场强加回路通量才是全部真实”的新图景。顺带一提，这也正是本章开头那个更朴素说法的量子升级版：绝对相位不可测（对应鸟脚下的五十万伏不可测），相位差可测（对应跨爪的电势差为零、跨翅的电势差致命），而相位差的账本是矢势记的。

这幅图景早已走出实验室，变成医院里每天都在用的技术。超导量子干涉仪（SQUID）就是一个把 AB 相位做成仪器的装置：一个超导环路的量子相位对穿过环的磁通敏感到以磁通量子为刻度，分辨率达到一个磁通量子的百万分之一量级。人脑神经元放电产生的磁场只有约 \(10^{-13}\) 特斯拉，是地磁场的十亿分之一，而 SQUID 阵列能隔着头骨把它测出来，绘成脑磁图，用来定位癫痫病灶、研究大脑如何响应声音和图像。当年“电子怎么会知道螺线管开了”的哲学困惑，如今是医生手里常规的检查单。
\end{mainline}

\step{脑洞实验与挑战题}

\begin{thought}[把地球充电到十亿伏]
做一个极端尺度的思想实验：假想某种装置把整个地球——连同大气、海洋和你——的静电势相对遥远的太空抬高了十亿伏。会发生什么吗？答案是：什么都不会发生。所有手机照常工作，所有鸟照常站在电线上，心电图照常跳动，因为每一个电路、每一个细胞感受的都只是自己内部的电势差，而所有电势被整体抬升同一个数，差值纹丝不动。这正是全局规范自由在行星尺度上的表演。但请再追问一步：为什么连“绝对电势”都定义不出来？因为想给绝对电势赋值，你需要把电荷搬到“无穷远”去对账，而宇宙不给你这个无穷远。再深一层：第 58 章说过，整体相位对称对应电荷守恒——也就是说，电荷既不会被创生也不会被消灭，连“给整个宇宙抬升零点”的原材料都凑不齐。描述的自由与守恒的硬约束，是同一枚硬币的两面。

再换一个极端尺度。宇宙学家设想过一种叫“宇宙弦”的假想天体：时空在大爆炸冷却后可能留下的细丝状伤疤，直径比原子核还小，绵延可达整个星系。如果宇宙弦存在，带电粒子绕着它飞一圈，即使弦外没有任何普通意义上的场，粒子的相位也可能被拧动一个固定的角度——整个星系尺度的阿哈罗诺夫—玻姆效应。天文学家真的分析过这类信号可能留下的观测痕迹。规范思想的力量在于：它从一只水瓶和一根高压线出发，却能一路问到宇宙最古老伤疤的门前。
\end{thought}

\begin{mainline}
最后，用一根指针把本章接到后面的章节。规范原理要求传播相互作用的粒子像光子一样没有质量——质量项会破坏逐点任选零点的自由。光子确实没有质量，电磁力力程无限，自洽；可是弱相互作用的 \(W\) 粒子质量高达约 80 吉电子伏特，换算成力程只有约 \(2\times10^{-18}\) 米，比一个质子还小一千倍，这就是为什么弱力在日常生活中完全隐身，只在原子核的衰变里露脸。一个没有质量的理论，怎么装下一个有质量的现实？直接给 \(W\) 写上质量项，理论会在高能下算出无穷大而崩溃；不写，又与实验矛盾。这道裂缝在 1960 年代逼出了二十世纪物理最漂亮的主意之一：不是方程自己破例，而是方程对称、世界所处的具体状态不对称。那是第 60 章的主角——对称性的自发破缺，以及给粒子“发质量”的希格斯场。本章从“描述不许依赖约定”出发，下一章将看到：有时候，世界为了守住方程的对称，偏偏挑了一个不对称的姿势落地。
\end{mainline}

\begin{puzzles}
\item 时区类比的体检。正文把矢势比作“时区换算表”，把局域相位自由比作“各市自定本地时间”。这个类比在哪里成立，在哪里失效？（思路提示：成立之处在于——各地的零点确实可以各选各的，而比较两地必须有换算规则，且物理结论不认零点。失效之处至少有两处：时区是离散的、全球协调好的约定，而规范零点在每一点是连续任意的；更要命的是，换算表本身只是纸面规则，而矢势沿回路的积分是可测量的物理——没有哪家航空公司的时刻表能让你的手表真的走时不同。好的类比送到一半就要学会告别。）
\item 为什么是 \(h/2e\)。实验室里把超导体制成一个环，发现穿过环的磁通只能取 \(h/2e\) 的整数倍，而不是本章磁通量子 \(h/e\) 的整数倍。系数里的那个 2 从哪来？这件事反过来告诉了我们超导体的什么秘密？（思路提示：绕环一圈，量子相位必须自洽——指针转过的总角度必须是整圈数，否则波函数在环上同一点有两个不同的值。相位积累公式里的 \(q\) 在超导体里不是电子电荷 \(e\)：超导体里的载流子是成对行动的电子（库珀对），\(q=2e\)，所以自洽条件给出 \(h/2e\)。这个不起眼的系数 2，正是“电子在超导体里两两结伴”的宏观铁证。）
\item 设计一个“测绝对电势”的实验。试试看，能不能设计出一个只用一个探头的仪器，真正测出“某一点的绝对电势”，而不悄悄引入第二个参考点？（思路提示：逐个排查你的方案——静电计的外壳、电压表的公共端、示波器的地线、乃至仪器操作员脚下的大地，每一个都在扮演“第二点”。你会发现所有方案要么偷偷接入了参考点，要么只能读出与仪器自身的差。这不是仪器不够精密，而是“绝对电势”在物理上根本不存在——它是描述层的约定，不是测量层的对象。）
\item 质量的麻烦。正文中说，规范原理要求传播粒子无质量，而实验上 \(W\) 粒子有质量，这个矛盾逼出了第 60 章的故事。请先用粗线条推理：为什么“传播粒子有质量”会导致力程变短？（思路提示：传递相互作用相当于在两点之间“借”出一个传播粒子。借出一个质量为 \(M\) 的粒子，意味着凭空凑出约 \(Mc^{2}\) 的能量账，量子理论允许这种账短暂存在，但借得越久越贵，于是交换只能在约 \(\hbar/(Mc)\) 的时间内完成，乘以光速就是力程。代入 \(M\) 约 80 吉电子伏特，力程约 \(2\times10^{-18}\) 米。光子质量为零，账可以无限期挂着，所以电磁力程无限。）
\end{puzzles}
